% Generated draft: requires mathematical and source review.
\documentclass[11pt]{article}
\usepackage[margin=25mm]{geometry}
\usepackage{fontspec}
\setmainfont{Latin Modern Roman}
\usepackage{amsmath,amssymb,mathtools,unicode-math}
\setmathfont{Latin Modern Math}
\usepackage{enumitem,needspace,xurl,hyperref}
\setlength{\parindent}{0pt}
\setlength{\parskip}{5pt}
\setlength{\emergencystretch}{3em}
\newcommand{\N}{\mathbb{N}}
\newcommand{\Z}{\mathbb{Z}}
\newcommand{\Q}{\mathbb{Q}}
\newcommand{\R}{\mathbb{R}}
\newcommand{\C}{\mathbb{C}}
\newcommand{\F}{\mathbb{F}}
\newcommand{\A}{\mathbb{A}}
\newcommand{\PP}{\mathbb{P}}
\newcommand{\E}{\mathbb{E}}
\newcommand{\eps}{\varepsilon}
\newcommand{\dd}{\,\mathrm{d}}
\begin{document}
\section*{UnsolvedMath ID: 20000858}
\textbf{Source problem code:} catalogue entry
\subsection{Definitions and mathematical statement}
This source is a talk-summary section, not a single formal conjecture. It reports Daniel Kersten's Bayesian influence-graph account of object perception, including discounting nuisance information such as lighting and cue integration. It reports perceptual organization associated with reduced V1 activity, increased LOC activity, and V1 prediction of percept on behavioral time scales. It explicitly leaves unresolved whether V1 reduction reflects predictive cancellation of well-predicted input or sparsification that enhances object features and suppresses others; the same aggregate fMRI observation was consistent with both. These are historical source reports, not new empirical findings.

The supplied mathematical partial result declares a common latent contrast \(\Delta x\in\R^n\) and a known linearized observation map \(y=H\Delta x\), \(H\in\R^{m\times n}\). Let \(\Delta x_P,\Delta x_S\) be two specified implementations of predictive cancellation and sparsification in that \emph{same} latent variable, and put \(\delta=\Delta x_P-\Delta x_S\). Then
\[
 H\Delta x_P=H\Delta x_S\quad\Longleftrightarrow\quad\delta\in\ker H.
\]
A scalar regional mean, \(H=n^{-1}(1,\ldots,1)\), leaves a kernel of dimension \(n-1\). For a fixed nonzero ambiguous \(\delta\), one added scalar row \(w^{\mathsf T}\) is minimal and sufficient to distinguish the pair exactly when \(w^{\mathsf T}\delta\neq0\). If \(\delta=0\), no additional linear measurement of that unchanged state distinguishes the implementations.

More generally, for a prespecified ambiguity subspace \(U\subset\ker H\) of dimension \(k\), added measurements \(W\Delta x\) identify all directions in \(U\) exactly when \(W|_U\) is injective, or \(\operatorname{rank}(W|_U)=k\). At least \(k\) scalar rows are required and \(k\) suffice if arbitrary suitable rows are available. Physical feasibility and independence of realizable measurements are not established by this linear-algebra count.

A same-mean example is \(\Delta x_P=(-2,0)\), \(\Delta x_S=(0,-2)\): both have mean \(-1\), whereas \(g=(1,-1)/\sqrt2\) gives \(g^{\mathsf T}\Delta x_P=-\sqrt2\), \(g^{\mathsf T}\Delta x_S=\sqrt2\). These are declared toy contrasts, not universal physiological predictions of the two mechanism names. The feature-tuning direction should be independently localized and fixed before the test. In an explicit Euclidean-error convention, if \(\|g\|=1\), \(\|x\|\leq M\), \(\|\widehat g-g\|\leq\epsilon_g\), and \(\|\widehat x-x\|\leq\epsilon_x\), then
\[
 |\widehat g^{\mathsf T}\widehat x-g^{\mathsf T}x|
 \leq\epsilon_x+M\epsilon_g+\epsilon_g\epsilon_x.
\]
An estimated sign is certified only if its magnitude exceeds that bound. The summary does not provide its original error notation or numerical calibration.

The aid proposes crossing organization \(o\in\{0,1\}\) with a feedback perturbation \(b\in\{0,1\}\). With condition means \(x_{bo}\), the feature-aligned difference-in-differences is
\[
 D=g^{\mathsf T}\bigl[(x_{11}-x_{10})-(x_{01}-x_{00})\bigr].
\]
A falsifiable experiment requires the declared mechanisms to make different predictions for this statistic and requires a justified causal interpretation of the feedback manipulation. The observation-rank theorem itself supplies neither intervention validity nor the biological answer.

The remaining historical discussion is also qualified. The notes say it was not yet known whether the described shape/color explaining-away cues were learned. Grossberg questioned priors in nonstationary environments and Bayesian accounts of Kanizsa illusions; Bruno disagreed and stressed choosing priors, while Dan mentioned a slow-motion prior. These are differing discussion positions, not proved impossibility or consistency theorems. Bill asked about discounting; Jeff asked about learning; Bill Softky asked about thalamic suppression, and Dan reported no such data then because fMRI localization of thalamus and especially LGN was difficult. Steve's Kanizsa illustration is mentioned without a further mathematical claim. None of these source-time limitations is converted into a current universal absence-of-data statement.
\subsection{Short English statement}
The talk summary leaves predictive cancellation and sparsification indistinguishable from aggregate V1 suppression. The supplied linearized model characterizes this ambiguity by a measurement null space and counts the extra scalar measurements needed to resolve a specified pair or subspace. Feature-sign margins and a proposed feedback experiment are conditional design tools, not a determination of the actual neural mechanism.
\subsection{Sources}
\begin{itemize}
\item Local registry: permanent UnsolvedMath ID \(20000858\), exact public code \texttt{catalogue entry}; \texttt{aim-biology-notes.json}, \texttt{notes[22]}. AIM domain Biology; registry category Miscellaneous; source tag section.
\item AIM 2003 workshop notes \emph{Inference and Prediction in Neocortical Circuits}, Chapter A, Daniel Kersten talk summary [AIM2003]: \url{https://aimath.org/WWN/brain/brain.pdf}. The verification field reports that the official PDF confirms the central passage; no independent PDF access is claimed here.
\item The section records discussion by Bill, Jeff, Grossberg, Bruno, Dan, Bill Softky, and Steve, with the qualifications reproduced above. Full participant identifications and underlying experiment citations are not supplied.
\item Supplied research summary, ``What aggregate V1 suppression cannot identify,'' attempt 1. No empirical dataset, calibrated observation operator, causal intervention protocol, or full supporting external bibliography is supplied; original remarks and literature fields are empty.
\end{itemize}
\end{document}
